\newpage
\section*{Đề bài}
\addcontentsline{toc}{section}{Đề bài}

\par Dựa vào lý thuyết ``Chapter 3 - Analysis of regression'' đã học, hãy tính toán 
và thực hiện các yêu cầu bên dưới.

\textbf{\underline{Yêu cầu:}}

\subsection*{Phần A: Dùng công cụ Excel, SPSS, Python và R tính toán.}

\begin{enumerate}

\item Giả sử bạn có dữ liệu trong Bảng 1 bên dưới và mức ý nghĩa là 5\%. 
Bạn đang nghiên cứu sự phụ thuộc của $Y$ với các biến độc lập $X_1$ và $X_2$ 
trong ngành F\&B.

\begin{table}[!htp]
\centering
\setlength{\tabcolsep}{0.05cm}
\renewcommand{\arraystretch}{1.2}

\resizebox{\linewidth}{!}{
\begin{tabular}{
|>{\centering\arraybackslash}p{0.2\linewidth}
|>{\centering\arraybackslash}p{0.2\linewidth}
|>{\centering\arraybackslash}p{0.2\linewidth}
|>{\centering\arraybackslash}p{0.2\linewidth}
|>{\centering\arraybackslash}p{0.2\linewidth}
|>{\centering\arraybackslash}p{0.2\linewidth}|
}
\hline
\textbf{Productivity (X1)} & \textbf{Produced Cost (X2)} & \textbf{Output of Product (Y)} 
& \textbf{Productivity (X1)} & \textbf{Produced Cost (X2)} & \textbf{Output of Product (Y)} \\
\hline
2.5 & 30 & 19.3 & 3.1 & 45 & 25.6 \\ \hline
3.4 & 30 & 22.4 & 3.2 & 40 & 24 \\ \hline
3.5 & 45 & 24.5 & 3.1 & 45 & 24.3 \\ \hline
3.6 & 40 & 26.9 & 3.5 & 45 & 27.3 \\ \hline
3.4 & 40 & 25.3 & 3.5 & 40 & 27.5 \\ \hline
3.3 & 45 & 22.7 & 3.4 & 30 & 21.2 \\ \hline
3.6 & 40 & 25.6 & 3.2 & 30 & 16.7 \\ \hline
3.5 & 40 & 27.2 & 2.5 & 30 & 18.4 \\ \hline
2.5 & 30 & 21.6 & 2.7 & 30 & 19.6 \\ \hline
\end{tabular}
}

\caption{Dữ liệu Productivity, Produced Cost và Output of Product}
\label{fig:table1}

\end{table}

\begin{enumerate}

\item[1.1] Dựa trên dữ liệu Bảng 1, hãy viết phương trình hồi quy bội 
(phương trình hồi quy tuyến tính) từ kết quả phân tích hồi quy và 
giải thích ý nghĩa của các chỉ số sau:

\begin{itemize}
\item Multiple R
\item R-square
\item Adjusted R-square
\item Significance F
\item Coefficient of independent variables
\item P-value
\end{itemize}

\item[1.2] Giả sử bạn có phương trình hồi quy bội (phương trình hồi quy phi tuyến tính) như sau:

\[
Y = b_0 + b_1 X_1 + b_2 X_2 + \frac{b_3}{X_1} + \frac{b_4}{X_1 X_2}
\]

Hãy viết phương trình hồi quy bội (phương trình hồi quy phi tuyến tính) 
từ kết quả phân tích thống kê. Dựa trên kết quả phân tích thống kê của bạn, 
hãy so sánh với kết quả phương trình hồi quy tuyến tính trong mục (1.1).

\end{enumerate}

\item Bạn có file dữ liệu \texttt{Home\_Market\_Value.xlsx} đính kèm, trong đó 
có 03 cột: \textbf{House Age}, \textbf{Square Feet}, và \textbf{Market Value}.

\begin{enumerate}

\item[2.1] Tạo phương trình hồi quy đơn biến với \textbf{Market Value} là 
biến phụ thuộc và \textbf{House Age} là biến độc lập. Hãy đánh giá và phân tích 
kết quả hồi quy của mô hình này.

\item[2.2] Tạo phương trình hồi quy đa biến với \textbf{Market Value} là biến 
phụ thuộc, \textbf{House Age} và \textbf{Square Feet} là hai biến độc lập. 
Hãy đánh giá và phân tích kết quả hồi quy của mô hình này.

\item[2.3] Hãy đánh giá và so sánh hai mô hình (2.1) và (2.2) thông qua các 
chỉ số sau:

\begin{itemize}
\item MAE
\item MAPE
\item MSE
\item RMSE
\item R-squared
\item Adjusted R-squared
\end{itemize}

\end{enumerate}

\end{enumerate}

\textbf{Chú ý:}

\begin{itemize}

\item Kết quả phần (A) phải được tóm tắt và chụp hình đưa vào tệp tin có quy cách 
đặt tên như sau:

\texttt{MSSV\_Họtên\_Lab03\_Output.docx} (\texttt{.pdf})

\item Gửi đính kèm các file Excel, SPSS, R và Python của Phần A với quy cách 
đặt tên như sau:

\begin{itemize}

\item \texttt{MSSV\_Họtên\_Lab03\_Excel.xlsx}
\item \texttt{MSSV\_Họtên\_Lab03\_SPSS.sav} (và \texttt{.spv})
\item \texttt{MSSV\_Họtên\_Lab03\_Python.ipynb}
\item \texttt{MSSV\_Họtên\_Lab03\_R.r}

\end{itemize}

\item Đóng gói lại thành 01 tập tin duy nhất và nộp theo quy cách đặt tên:

\texttt{MSSV\_Họtên\_Lab03.rar} (\texttt{.zip})

\item Deadline nộp bài: \textbf{24/03/2026}.

\item Sinh viên nộp bài không đúng deadline sẽ không được tính điểm Lab03.

\item Sinh viên nộp không đủ file hoặc đặt tên sai quy cách sẽ bị trừ điểm.

\item Các bài có hình thức trình bày và nội dung giống nhau sẽ nhận điểm \textbf{Zero}.

\end{itemize}